\documentclass[11pt]{article}

% -------------------------------------------------
% Packages
% -------------------------------------------------
\usepackage[utf8]{inputenc}
\usepackage[T1]{fontenc}
\usepackage{lmodern}

\usepackage{amsmath, amssymb, amsthm, mathtools}
\usepackage{bm}

\usepackage{graphicx}
\usepackage{subcaption}
\usepackage{booktabs}
\usepackage{multirow}

\usepackage{geometry}
\geometry{margin=1in}

\usepackage{hyperref}
\hypersetup{
    colorlinks=true,
    linkcolor=blue,
    citecolor=blue,
    urlcolor=blue
}

\usepackage{enumitem}
\usepackage{float}

\usepackage{authblk}

% -------------------------------------------------
% Theorem environments
% -------------------------------------------------
\newtheorem{theorem}{Theorem}
\newtheorem{proposition}{Proposition}
\newtheorem{remark}{Remark}

% -------------------------------------------------
% Commands
% -------------------------------------------------
\newcommand{\R}{\mathbb{R}}
\newcommand{\E}{\mathbb{E}}
\newcommand{\norm}[1]{\left\lVert #1 \right\rVert}
\newcommand{\ip}[2]{\left\langle #1, #2 \right\rangle}

% -------------------------------------------------
% Title
% -------------------------------------------------
\title{
A Small Empirical Study of the Neural Tangent Kernel Regime
}

\author{
Desbois-Renaudin Méo
}

\affil{
Independent Study
}

\date{\today}

% -------------------------------------------------
% Document
% -------------------------------------------------
\begin{document}

\maketitle

% -------------------------------------------------
% Abstract
% -------------------------------------------------
\begin{abstract}

The Neural Tangent Kernel (NTK) framework predicts that sufficiently wide neural networks trained with gradient descent behave approximately as kernel methods.
In this note, we explore how this ``lazy training'' regime emerges in finite-width multilayer perceptrons trained on the handwritten digits dataset.
We empirically investigate several quantities associated with the NTK regime, including parameter drift, stability of the empirical NTK during training, and dependence on network width.
Our goal is not to establish new theoretical results, but rather to provide a compact and pedagogical exploration of finite-width behavior through the lens of NTK theory.

\end{abstract}

% -------------------------------------------------
\section{Introduction}

The Neural Tangent Kernel (NTK) paradigm introduced by Jacot et al.~\cite{jacot2018ntk} provides a striking perspective on the training dynamics of overparameterized neural networks.
In the infinite-width limit, gradient descent training can be approximated by kernel regression with a deterministic kernel.

This framework has generated significant interest in deep learning theory, especially through later generalizations to broad neural network architectures.
However, practical neural networks are finite, and it is natural to ask to what extent finite-width networks actually exhibit lazy training behavior.

In this note, we perform a small empirical investigation of this question on simple multilayer perceptrons trained on MNIST dataset.
We focus on experimentally measurable quantities associated with the NTK regime and study their evolution as network width increases.

% -------------------------------------------------
\section{Background}

% -------------------------------------------------
\subsection{Neural Tangent Kernel}

Consider a neural network defined with its parameters $\theta \in \R^p$ by the function : 
\[
f(.; \theta): \R^{n_0} \rightarrow \R^{n_L}
\]

The Neural Tangent Kernel associated with the network is defined as
\[
\Theta(x, x')
=
\nabla_\theta f(x; \theta)^T
\nabla_\theta f(x'; \theta)
\]

In the infinite-width limit, under suitable initialization and scaling assumptions, the NTK converges to a deterministic kernel and remains approximately constant during training.

% -------------------------------------------------
\subsection{Lazy Training Regime}

The lazy training regime refers to a setting where:
\begin{itemize}
    \item parameters move little during training,
    \item the NTK remains nearly constant,
    \item the network behaves approximately as its first-order linearization around initialization.
\end{itemize}

A central question is how strongly finite-width networks exhibit these behaviors in practice.

% -------------------------------------------------
\section{Experimental Setup}

% -------------------------------------------------
\subsection{Dataset}

We use the handwritten digits dataset from scikit-learn.
The dataset contains grayscale images of handwritten digits represented as $8 \times 8$ images flattened into vectors of dimension $64$.

We standardize the input features using feature-wise normalization.

% -------------------------------------------------
\subsection{Architecture}

We consider a two-layer multilayer perceptron of the form
\[
f(x)
=
W_2 \sigma(W_1 x + b_1) + b_2,
\]
where $\sigma$ is the ReLU activation function.

We study several hidden widths:
\[
m \in \{10, 50, 100, 500, 1000\}.
\]

% -------------------------------------------------
\subsection{Optimization}

The models are trained using gradient descent (or Adam).
For each width, we repeat training across multiple random initializations.

% -------------------------------------------------
\subsection{Measured Quantities}

We study the following quantities:

\paragraph{Parameter drift}
\[
\norm{\theta_t - \theta_0}.
\]

\paragraph{Kernel drift}
\[
\norm{K_t - K_0}_F,
\]
where $K_t$ denotes the empirical NTK at training step $t$.

\paragraph{Test accuracy}

We also track standard predictive performance.

% -------------------------------------------------
\section{Results}

% -------------------------------------------------
\subsection{Width and Predictive Stability}

% Example figure placeholder
\begin{figure}[H]
    \centering
    \includegraphics[width=\linewidth]{figures/comp_widths_L4_v2.png}
    \caption{
    Test accuracy as a function of network width.
    Shaded regions indicate one standard deviation across random initializations.
    }
    \label{fig:accuracy}
\end{figure}

As network width increases, training becomes increasingly stable across random initializations.

% -------------------------------------------------
\section{Discussion}

Our experiments suggest that finite-width multilayer perceptrons progressively approach the lazy training regime as width increases.
In particular, both parameter drift and NTK drift decrease with width.

At the same time, finite networks remain meaningfully different from their infinite-width idealization.
This highlights the importance of understanding when NTK approximations are accurate in practice.

The present note is intentionally modest in scope.
Our goal is primarily pedagogical and exploratory rather than theoretical.

% -------------------------------------------------
\section{Conclusion}

We presented a small empirical exploration of finite-width behavior through the lens of the Neural Tangent Kernel framework.
Even in simple settings, increasing width leads to increasingly stable optimization dynamics consistent with lazy training predictions.

Future directions include studying convolutional architectures, explicit linearization comparisons, and departures from the NTK regime in feature-learning settings.

% -------------------------------------------------
\bibliographystyle{plain}

\begin{thebibliography}{9}

\bibitem{jacot2018ntk}
Arthur Jacot, Franck Gabriel, and Clément Hongler.
\newblock Neural Tangent Kernel: Convergence and Generalization in Neural Networks.
\newblock In \emph{NeurIPS}, 2018.

\bibitem{yang2020tensor}
Greg Yang.
\newblock Tensor Programs II: Neural Tangent Kernel for Any Architecture.
\newblock arXiv preprint arXiv:2006.14548, 2020.

\end{thebibliography}

\end{document}